\section{The case of a ground field}
\label{II.5}
% original p. 52

\begin{proposition}
\label{II.5.1}
Let $k$ be a field, $X$ a prescheme of finite type over~$k$,
$x$ a point of~$X$ and $n$ the dimension of~$X$ at~$x$,
$f\colon X\to\Spec k[t_1,\dots,t_n]=Y$
a morphism, defined by elements
$f_i\in\Gamma(X,\cal{O}_X)$.
The following conditions are equivalent (and imply that
$X$ is smooth over~$k$ at~$x$, and a fortiori regular at~$x$
by~\ref{II.3.1}):
\begin{enumerate}
\item[(i)] $f$ is \'etale at~$x$.
\item[(ii)] The $df_i$ form a basis of~$\it{\Omega}_{X/k}^1$ at~$x$.
\item[(iii)] The $df_i$ generate $\it{\Omega}_{X/k}^1$ at~$x$.
\end{enumerate}
\end{proposition}

Since (i) implies that $X$ is smooth over~$k$ at~$x$, the implication
(i)$\To$(ii) is a special case of~\ref{II.4.8},
(ii)$\To$(iii) is trivial, it remains to prove
(iii)$\To$(i). Now if (iii) is satisfied, $f$ is net
at~$x$ by virtue of lemma~\ref{II.4.1}, hence (replacing $x$ by
an open neighborhood) quasi-finite, hence dominant for reasons of
dimensions. Since $Y$ is regular, it follows that $f$ is
\'etale by (I~\SourceRef{I.9.5}{9.5}~(ii)) or (I~\SourceRef{I.9.11}{9.11}).

\begin{corollary}
\label{II.5.2}
Under the preliminary conditions of~\ref{II.5.1}, suppose that
$\kres(x)$
is a \emph{finite separable} extension of~$k$, and
that the $f_i$ ($1\leq i\leq n$) define elements
of~$\goth{m}_x$. Then the preceding conditions
are equivalent to
\begin{enumerate}
\item[(iv)] The $f_i$ form a system of generators of
$\goth{m}_x$ (or equivalently: the $f_i$ mod $\goth{m}_x^2$ form a basis
of~$\goth{m}_x/\goth{m}_x^2$
over~$\kres(x)$).
\end{enumerate}
\end{corollary}

Indeed, (iv)$\To$(iii) by virtue of the exact sequence
\begin{equation}
\label{eq:II.5.1}
\goth{m}_x/\goth{m}_x^2\to\Omega_{\cal{O}_x/k}^1\to
\Omega_{\kres(x)/k}^1\to 0
\end{equation}
and taking into account that~$\Omega_{\kres(x)/k}^1=0$ since $\kres(x)$ is
\'etale over~$k$. On the other hand (ii) implies (iv), for since $X$ and
$\Spec(\kres(x))$ are smooth over~$k$ at~$x$, one can in the preceding
exact sequence put a $0$ on the left by virtue of~\ref{II.4.10}~(iv).

\begin{corollary}
\label{II.5.3}
Let $x$ be a point of~$X$ (of finite type over~$k$). If $X$ is
smooth over~$k$ at~$x$, then $\cal{O}_x$ is regular, the converse
being true if
$\kres(x)$
is a finite separable
extension of~$k$.
\end{corollary}

Indeed,
% original p. 53
the converse follows from~\ref{II.5.2}, by taking a
regular system $(f_i)$ of generators of
$\goth{m}_x$. (N.B. instead of~\ref{II.5.2}, one can also invoke
theorem~\ref{II.4.15}). One concludes:

\begin{proposition}
\label{II.5.4}
Let $X$ be a prescheme of finite type over~$k$. If $X$ is smooth
over~$k$, it is regular, the converse being true if $k$
is perfect.
\end{proposition}

For the converse, one notes that by virtue of~\ref{II.5.3}, $X$ is
smooth over~$k$ at every closed point, hence everywhere since the set
of points where it is smooth is open.

\begin{theorem}
\label{II.5.5}
Let $X$ be a prescheme of finite type over~$k$, $x$ a point of
$X$, $n$ the dimension of~$X$ at~$x$, $k'$ a perfect extension of
$k$. The following conditions are equivalent:
\begin{enumerate}
\item[(i)] $X$ is smooth over~$k$ at~$x$.
\item[(ii)] $\it{\Omega}_{X/k}^1$ is free of rank $n$ at~$x$.
\item[(ii~bis)] $\it{\Omega}_{X/k}^1$ is generated by $n$
elements at~$x$.
\item[(iii)] $X$ is differentially smooth over~$k$ at~$x$.
\item[(iv)] There exists an open neighborhood $U$ of~$x$ such that
$U\otimes_k k'$ is regular (\ie the local rings of its
points are regular).
\end{enumerate}
\end{theorem}

One has (i)$\To$(ii) by~\ref{II.4.3}~(ii), (ii)$\To$(ii~bis) trivially and (ii~bis)$\To$(i) thanks
to~\ref{II.5.1}. Since $X$ is flat over~$k$, one has
(i)$\Leftrightarrow$(iii) by virtue of~\ref{II.4.18}. One has
(i)$\To$(iv) since smooth is invariant under extension of the
base and implies regular, and (iv)$\To$(i) for by
Proposition~\ref{II.5.4}, one sees that $U\otimes_k k'$ is simple over~$k'$, hence $U$ is simple over~$k$ by~\ref{II.4.13}.

Taking for $x$ the generic point of~$X$ assumed
irreducible, one finds:

\begin{corollary}
\label{II.5.6}
Let $K$ be a local Artin ring which is a localization of an algebra of
finite type over the field~$k$ (for example, $K$ may be an
extension of finite type of~$k$), let~$n$ be the transcendence
degree over~$K$ of its residue field. Equivalent conditions:
\begin{enumerate}
\item[(i)] $K$ is a finite separable extension of a purely
transcendental extension $k(t_1,\dots,t_n)$ of~$k$.
\item[(ii)]
% original p. 54
$\it{\Omega}^1_{X/k}$
is a $K$-module free of rank~$n$.
\item[(ii~bis)]
$\it{\Omega}^1_{X/k}$
is a $K$-module admitting $n$
generators.
\item[(iii)] The completion $O'$ of~$K\otimes_k K$ for the
topology defined by the powers of the augmentation
ideal $K\otimes_k K\to K$ is a
``regular'' augmented
$K$-algebra,
\ie isomorphic to an
algebra of formal series in~$K$ (N.B. If~$K$ is a field,
this is equivalent to saying that~$O'$ is a
regular local ring).
\item[(iv)] $K$ is a separable extension of~$k$.
\end{enumerate}
\end{corollary}

Indeed, one can always consider $K$ as the local ring of the
generic point of an irreducible scheme of finite type
$X$ over~$k$, and the conditions envisaged are
the conditions of the same name in~\ref{II.5.5}, taking in~(iv)
for~$k'$ an algebraically closed extension of~$k$. Only
the implication $K$ separable over~$k$ $\To$ $X$ smooth over~$k$ at~$x$,
requires a proof. Now one is immediately
reduced thanks to~\ref{II.4.13} to the case where the ground
field is~$k'$, hence algebraically closed, hence where there exists a
point $a$ of~$X$ rational over~$k$. But then $X$ is smooth over~$k$
at~$a$ by~\ref{II.5.4}, a fortiori it is smooth over~$k$ at the
generic point~$x$, qed\footnote{\Cf Errata at the end of the present Exp\ptbl II
(p\ptbl \pageref{II.fin.errata})}.

One will note that in the case where~$K$ is an extension of finite
type of~$k$, the equivalence of
(i), (ii), (ii~bis), (iv)
is well
known, but that we have not made use of any of
these
already known
equivalences.
Of course,
Proposition~\ref{II.5.1} contains as a special case that a sequence
of elements~$x_i$ ($1\leq i\leq n$) is a ``separating
transcendence basis'' of~$K$ over~$k$ if and only if
the~$dx_i$ form a basis of the $K$-module $\Omega^1_{K/k}$, a fact well
known moreover.

\begin{corollary}
\label{II.5.7}
Let $X$ be a prescheme of finite type over a
field~$k$. For $X$ to be smooth over~$k$, it is necessary and sufficient that
$\it{\Omega}^1_{X/k}$
be locally free, and that the local rings of the
generic points of the irreducible components of~$X$
be separable extensions of~$k$ (this last
condition being automatically satisfied if $k$ is
perfect and $X$ reduced).
\end{corollary}

One may assume $X$ connected, let $n$ be the rank of
$\it{\Omega}^1_{X/k}$
assumed locally free. By the hypothesis
and~\ref{II.5.6}, this is also the transcendence degree of the
extensions
% original p. 55
of~$k$ defined by the local rings of the generic
points of~$X$, hence all the
irreducible components of~$X$ are of dimension~$n$. One then concludes
thanks to~\ref{II.5.5}.

One should take care that if $K$ is a finite extension (not
necessarily separable) of~$k$, then $\Omega^1_{K/k}$ is
a free $k$-module, hence setting $X=\Spec(K)$, $\Omega^1_{X/k}$ is a
locally free sheaf, and $X$ is reduced, without $X$ being
necessarily smooth over~$k$. Extending then the scalars to
the algebraic closure of~$k$, one finds an analogous example,
where $k$ is algebraically closed, but $X$ on the other hand
not being reduced.

\begin{corollary}
\label{II.5.8}
Let $X$ be a prescheme of finite type over the field~$k$, $x$ a
point of~$X$, $n$ the dimension of~$X$ at~$x$, $p$ the dimension
of~$\cal{O}_x$,
\ie the codimension in $X$ of the closure $Y$ of
$x$ in~$X$; hence $n-p$ the transcendence degree
of~$\kres(x)$
over~$k$. Let $f_i$ ($1\leq i\leq n$) be elements
of~$\cal{O}_x$, such that $f_i\in \goth{m}_x$ for $1\leq i\leq p$. The
following conditions are equivalent

\begin{enumerate}
\item[(i)] the germ of morphism at~$x$
\[
X\to \Spec\bigl(k[t_1,\dots,t_n]\bigr)
\]
defined by the $f_i$ is \'etale at~$x$.
\item[(ii)] The $f_i$ ($1\leq i\leq p$)
generate~$\goth{m}_x$,
\ie form a regular system of parameters
of~$\cal{O}_x$, and the classes in
$\kres(x)$ of the~$f_j$ ($p+1\leq j\leq n$) form a separating
transcendence basis, (\ie the~$d\overline{f}_j$ ($p+1\leq j\leq n$) form a basis
of~$\Omega^1_{\kres(x)/k}$, or equivalently generate~$\Omega^1_{\kres(x)/k}$).
\end{enumerate}
\end{corollary}

Suppose~(i) is satisfied. It follows that the $df_i(x)$
form a basis of~$\it{\Omega}^1_{X/k}(x)$ \eqref{II.4.8} hence their
images $d\overline{f}_i(x)$ in $\Omega^1_{\kres(x)/k}$ generate
this vector space over~$k$. Since the $\overline{f}_i$ for $1\leq
i\leq p$ are zero, it follows that it suffices to take the
$d\overline{f}_i(x)$ with $p+1\leq i\leq n$. Since the
transcendence degree of~$\kres(x)$ over~$k$ is~$n-p$, it then follows
from corollary~\ref{II.5.6} criterion~(iii) (applied
to~$K=\kres(x)$) that $Y$ is smooth over~$k$ at its
generic point~$x$, and that the~$d\overline{f}_i(x)$ ($p+1\leq i\leq
n$) form a \emph{basis} of~$\Omega^1_{\kres(x)/k}$
over~$\kres(x)$. Consequently, condition~(ii) of~\ref{II.4.9} is
satisfied, hence also condition~(iii) and in particular
the~$f_i$ ($1\leq i\leq p$) form a system of
generators of~$\goth{m}_x$. Since
% original p. 56
$\cal{O}_x$ is of dimension~$p$, they therefore form a
regular system of parameters at~$x$. This proves~(ii).

Suppose~(ii) is satisfied. By virtue of the exact
sequence~\eqref{II.5.1}, it follows that the~$df_i(x)$
generate~$\it{\Omega}^1_{X/k}$,
whence~(i) thanks to prop\ptbl\ref{II.5.1}.

\begin{corollary}
\label{II.5.9}
Let~$X$ be a prescheme of finite type over the
field~$k$, $x$ a point of~$X$, $n$ the dimension of~$X$ at~$x$, $p$ the
dimension
of~$\cal{O}_x$,
\ie the codimension of the closure~$Y$
of~$x$ in~$X$, hence~$n-p$ the transcendence degree
of~$\kres(x)$
over~$k$. Equivalent conditions:
\begin{enumerate}
\item[(i)] $\cal{O}_x$ is regular
and~$\kres(x)$
 is a
separable extension of~$k$.
\item[(ii)] $X$ is smooth over~$k$ at~$x$, and the canonical homomorphism
\[
\goth{m}_x/\goth{m}_x^2\to
\Omega^1_{\cal{O}_x/k}\otimes_{\cal{O}_x}
\kres(x)=\it{\Omega}^1_{X/k}(x)
\]
is injective.
\item[(iii)] There are~$f_i\in \cal{O}_x$ ($1\leq i\leq n$)
with~$f_i\in \goth{m}_x$ for~$1\leq i\leq p$, such that the germ of
morphism at~$x$ from~$X$ to~$\Spec\bigl(k[t_1,\dots,t_n]\bigr)$
defined by the~$f_i$ is \'etale at~$x$ (\ie by~\ref{II.5.1}
such that the~$df_i(x)$ generate~$\it{\Omega}^1_{X/k}(x)$).
\item[(iv)] There are~$f_i\in \cal{O}_x$ ($1\leq i\leq n$) such that
the~$f_i$ ($1\leq i\leq p$) generate~$\goth{m}_x$ and that
the~$df_j(x)$ ($p+1\leq j\leq n$)
generate~$\Omega^1_{\kres(x)/k}$ over~$\kres(x)$.
\end{enumerate}
\end{corollary}

The equivalence of~(iii) and~(iv) follows from
corollary~\ref{II.5.8}, these conditions are also equivalent
by~\ref{II.4.9} to the fact that~$X$ is smooth over~$k$ at~$x$ and
that condition~(ii)
of~\ref{II.4.10}
is satisfied. Hence they are equivalent to the fact that~$X$ is
smooth over~$k$ at~$x$ and that condition~(iv) of~\ref{II.4.10} is
satisfied, hence to~\ref{II.5.9}~(ii). Or to the fact that~$X$ is
smooth over~$k$ at~$x$ and that condition~(i) of~\ref{II.4.10} is
satisfied, which here simply means
that~$\kres(x)$
is
separable over~$k$. This implies~\ref{II.5.9}~(i), it remains to
prove that~\ref{II.5.9}~(i) implies it, \ie to prove the

\begin{corollary}
\label{II.5.10}
Let~$x$ be a point
of a
prescheme of finite type over the field~$k$, such
that~$\kres(x)$
is separable over~$k$. For~$X$ to be smooth over~$k$ at~$x$,
it is necessary and sufficient that it be regular at~$x$ (\ie that the local
ring~$\cal{O}_x$ be regular).
\end{corollary}

Indeed, if this is so, one can evidently find
$f_i\in \cal{O}_x$ ($1\leq i\leq n$) satisfying
condition~\ref{II.5.9}~(iv).

\subsection*{Errata}
\label{II.fin.errata}
In
% original p. 57
the present no., proof of~\ref{II.5.6}, one
has used the fact that a nonempty reduced scheme of finite type
over an algebraically closed field admits at least one
regular point (hence smooth), a fact which is usually proved by
the differential method (via Zariski's theorem that
the set of regular points of~$X$ is open). If one wants
to avoid a vicious circle, one must prove that if~$K/k$ is
a separable extension of finite type, and if the~$f_i\in K$ are
such that the~$d_{K/k}f_i$ form a basis of~$\Omega^1_{K/k}$, ($1\leq
i\leq n$), then~$n$ is the transcendence degree of~$K$
over~$k$,
\ie the~$f_i$ are algebraically independent. The
proof of this fact with the aid of the
Mac\kern2pt Lane
criterion is well known, \cf Bourbaki, Alg\`ebre, Chap\ptbl V par\ptbl 9 th\ptbl 2: one
takes a polynomial~$g\in k[t_1,\dots,t_n]$ of minimal
degree such that~$g(f_1,\dots,f_n)=0$, one has therefore
\[
\sum\frac{dg}{d t_i}(f_1,\dots,f_n) df_i=0
\]
whence (since the~$df_i$ form a basis of~$\Omega^1_{K/k}$) the
fact that the~$dg/d t_i$ vanish at~$(f_1,\dots,f_n)$, hence are zero
by the minimal character of~$g$, hence if~$k$ is of
characteristic~$0$ one has~$g=0$, and if~$k$ is of
characteristic~$p\neq 0$ one has~$g=h(t_1^p,\dots,t_n^p)$.
Using the
Mac\kern2pt Lane
criterion,
one sees that the
polynomial~$h\in k[t_1,\dots,t_n]$ also vanishes at~$(f_1,\dots,f_n)$,
whence again~$g=0$ by the minimal character of~$g$.
